% TOP_PROBLEM: {"id": 1101280, "title": "Questions in Geometric Group Theory — Q 12.80", "category_id": 17, "category": {"id": 17, "name": "group_theory", "display_name": "Group Theory", "description": "Problems about groups, group actions, representations, and related algebraic structures.", "slug": "group-theory", "order_index": 17, "created_at": "2026-07-31T15:26:25.671Z"}, "status": "open", "problem_number": "AMR-010-1280", "set_id": 13, "statusQualification": "States the standard countable-discrete Haagerup formulation, including the finitely generated case relevant to the source."}
% FIRST_POSTED: 2026-10-01
% ENTRY_KIND: statement_only
% UnsolvedMath ID 1101280; source catalogue code AMR-010-1280
% !TeX program = lualatex
% Definition and sources only; this file is not an LLM solution attempt.
\documentclass[11pt,a4paper]{article}
\usepackage[margin=25mm]{geometry}
\usepackage{fontspec}
\setmainfont{lmroman10-regular.otf}[BoldFont=lmroman10-bold.otf,
 ItalicFont=lmroman10-italic.otf,BoldItalicFont=lmroman10-bolditalic.otf]
\usepackage{amsmath,amssymb,mathtools}
\usepackage{unicode-math}
\setmathfont{latinmodern-math.otf}
\AtBeginDocument{\let\setminus\smallsetminus}
\usepackage{xurl,hyperref}
\hypersetup{colorlinks=true,urlcolor=blue}
\setcounter{secnumdepth}{0}
\setlength{\parindent}{0pt}
\setlength{\parskip}{5pt}
\setlength{\emergencystretch}{3em}
\begin{document}
\section*{Questions in Geometric Group Theory — Q 12.80}\label{1101280}
{\small UnsolvedMath ID 1101280 · AMR-010-1280 · Group Theory · AMR Open Problem Lists}
\subsection{Definitions and mathematical statement}
A group $N$ is nilpotent if its lower central series, defined by $\gamma_1N=N$ and $\gamma_{j+1}N=[N,\gamma_jN]$, reaches the trivial subgroup after finitely many steps. It is torsion-free if no nonidentity element has finite order. A group $G$ is residually torsion-free nilpotent if, for every $g\ne1$, there is a homomorphism $\phi:G\to N$ to some torsion-free nilpotent group with $\phi(g)\ne1$. The nilpotency class and quotient may depend on $g$.

For a countable discrete group, the Haagerup property means existence of a metrically proper affine isometric action on a real Hilbert space $\mathcal H$. Such an action has form
\[
 g\cdot v=\pi(g)v+b(g),\qquad
 b(gh)=b(g)+\pi(g)b(h),
\]
where $\pi$ is an orthogonal representation. Properness means that $\{g:\|b(g)\|\leq R\}$ is finite for every finite $R$.

Must every countable residually torsion-free nilpotent group have the Haagerup property? In particular, ask this for finitely generated groups, the natural setting of the source list. The proper Hilbert action need not factor through any single nilpotent quotient. Residual separation supplies many quotient maps, whereas the conclusion asks for one action whose displacement eventually escapes every bounded set. No common nilpotency class, uniform residual bound, or finite-dimensional Hilbert space is assumed. Mere absence of nontrivial torsion is weaker than the stated residual hypothesis.
\subsection{Short English statement}
Does residual torsion-free nilpotence guarantee a proper affine isometric action on Hilbert space?
\subsection{Sources}
\begin{itemize}
\item[S1] ULAM.AI and the UnsolvedMath Contributors, supplied problems.json, record 1101280 (AMR-010-1280), AMR Open Problem Lists. Catalogue identity and source transcription; CC BY 4.0.
\url{https://huggingface.co/datasets/ulamai/UnsolvedMath}.
\item[S2] Bestvina - Questions in Geometric Group Theory (pdf) (2004), Question 12.80 (PDF page 21). Original source indexed by AMR-010-1280.
\url{https://www.math.utah.edu/~bestvina/eprints/questions-updated.pdf}.
\end{itemize}
\end{document}
